% TOP_PROBLEM: {"id": 4700031, "title": "Conjecture of multiplicative persistence", "category_id": 2, "category": {"id": 2, "name": "combinatorics", "display_name": "Combinatorics", "description": "Counting problems, graph theory, discrete structures.", "slug": "combinatorics", "order_index": 2, "created_at": "2026-07-31T15:26:25.671Z"}, "status": "partially_solved", "problem_number": "AMR-046-0031", "set_id": 13}
% FIRST_POSTED: 2026-10-02
% ENTRY_KIND: statement_only
% UnsolvedMath ID 4700031; source catalogue code AMR-046-0031
% !TeX program = lualatex
% Definition and sources only; this file is not an LLM solution attempt.
\documentclass[11pt,a4paper]{article}
\usepackage[margin=25mm]{geometry}
\usepackage{fontspec}
\setmainfont{lmroman10-regular.otf}[BoldFont=lmroman10-bold.otf,
 ItalicFont=lmroman10-italic.otf,BoldItalicFont=lmroman10-bolditalic.otf]
\usepackage{amsmath,amssymb,mathtools}
\usepackage{unicode-math}
\setmathfont{latinmodern-math.otf}
\AtBeginDocument{\let\setminus\smallsetminus}
\usepackage{xurl,hyperref}
\hypersetup{colorlinks=true,urlcolor=blue}
\setcounter{secnumdepth}{0}
\setlength{\parindent}{0pt}
\setlength{\parskip}{5pt}
\setlength{\emergencystretch}{3em}
\begin{document}
\section*{Conjecture of multiplicative persistence}\label{4700031}
{\small UnsolvedMath ID 4700031 · AMR-046-0031 · Combinatorics · AMR Open Problem Lists}
\subsection{Definitions and mathematical statement}
Write a positive integer $n$ in decimal notation without leading zeros,
$n=\sum_{j=0}^{r-1}d_j10^j$, with $d_j\in\{0,\ldots,9\}$ and
$d_{r-1}\ne0$. Define $\Pi(n)=\prod_{j=0}^{r-1}d_j$ and put
$\Pi(0)=0$. Every one-digit nonnegative integer is fixed by $\Pi$.
Define its usual multiplicative persistence by
\[
 P(n)=\min\{k\ge0:\Pi^k(n)\in\{0,1,\ldots,9\}\}.
\]
The minimum exists because the digit product of a multi-digit positive
integer is smaller than the integer. The conjecture is
\[
 P(n)\le11\qquad\text{for every positive integer }n.
\]
Thus no decimal integer should require twelve successive digit-product
operations before reaching a one-digit fixed point.

The catalogue uses the least positive $k$ with
$\Pi^k(n)=\Pi^{k+1}(n)$. This equals $\max\{1,P(n)\}$, so it gives
exactly the same upper-bound question. This convention matters only for
integers that already have one digit.

The base ten is fixed; changing the base produces a different persistence
problem. Zeros in a multi-digit integer immediately give product zero,
while the order of its digits does not affect its first product. These
reductions help organize searches but do not bound all possible input
lengths. A finite computation, however extensive, does not prove the
universal bound without a theorem reducing arbitrary decimal lengths
to the checked range.
\subsection{Short English statement}
Must repeated multiplication of decimal digits reach a one-digit number within eleven steps?
\subsection{Sources}
\begin{itemize}
\item[S1] ULAM.AI and the UnsolvedMath Contributors, supplied problems.json, record 4700031 (AMR-046-0031), AMR Open Problem Lists. Catalogue identity and source transcription; CC BY 4.0.
\url{https://huggingface.co/datasets/ulamai/UnsolvedMath}.
\item[S2] Gasull - Some open problems in low dimensional dynamical systems (2020), Problem environment 31: Conjecture of multiplicative persistence. Original source indexed by AMR-046-0031.
\url{https://arxiv.org/abs/2012.02524}.
\end{itemize}
\end{document}
